%Abstract
\setlength{\headheight}{14pt}
\begin{center}
	{\huge \bf Abstract}
	\line(1,0){430}
\end{center}

%The abstract should contain a summary of the work presented in the report in a single paragraph. The paragraph should provide a general background of the problem, methodology and results.

Large language models now answer hard scientific questions, yet they still hallucinate and follow broken lines of reasoning. Multi-agent debate offers one way to push back on this. When several models argue a question out before committing to an answer, the reasoning usually gets sharper and errors get caught.

The weakness is in how these debates settle. Most systems pick a winner by majority vote or by which agent sounds most confident. When a confident majority is simply wrong, it can pull a correct minority agent off its answer, and the group lands on the wrong conclusion. No current method checks whether an agent's claims actually hold up against outside scientific evidence before letting that agent sway the room.

The framework addresses this problem by adjusting how much each agent counts as
the debate unfolds.

The framework breaks every agent response into small factual claims. For each claim it retrieves relevant scientific literature and checks whether the evidence backs the claim or contradicts it. Every agent then carries a trust score that moves up or down as these checks accumulate over the rounds of debate. The final answer is built from trust-weighted aggregation rather than a raw count of votes.

The trust update assigns less influence to unsupported claims, allowing a correct
minority agent to retain influence even when outnumbered.

The framework is expected to reduce hallucinations and keep reasoning tied to
verifiable sources. It may also reduce premature agreement on a wrong answer while
operating without model fine-tuning.

\clearpage
\setlength{\headheight}{12pt}
